\documentclass[a4paper,12pt]{article}
\usepackage[utf8]{inputenc}
\usepackage{amsmath}
\usepackage{amssymb}
\usepackage{graphicx}
\usepackage{array}
\usepackage{xcolor}
\usepackage{geometry}
\usepackage{longtable}
\usepackage{multicol}
\usepackage{booktabs}
\usepackage{verbatim}
\usepackage{hyperref}

\geometry{margin=2.5cm}

\begin{document}

%Hier noch ein Titel gemäß Nebenfach/Lehramt Ia

\section*{Glossar für Python-Befehle}

\begin{center}
\begin{tabular}{|p{3cm}|p{8cm}|p{3cm}|}
\hline
\textbf{Befehl} & \textbf{Beschreibung} & \textbf{Beispiel} \\
\hline
\end{tabular}
\end{center}

%Villeicht lieber weniger Abstand mit einer Tabelle und zwei \hline

\begin{center}
\begin{tabular}{|p{3cm}|p{8cm}|p{3cm}|}
\hline
a = b & a := b Der Wert von a wird auf b gesetzt. Die Leerzeichen sind redundant und nur für Lesbarkeit & a = 2 \\
\hline
print(x) & Gibt x aus & print(a) \newline{} \texttt{=>} 2\\
\hline
print(x, y, "z") & Innerhalb eines prints können mehrere Variablen verkettet werden& print("a = " , a) \newline{} \texttt{=>} a = 2\\
\hline
import \textcolor{gray}{Package} as \textcolor{gray}{Abkürzung} & Läd Bibliotheken m ihre Befehle nutzen zu können. Das "as Abkürzung" ist hierbei nicht erforderlich, sondern verringert die Schreibarbeit. & import numpy as np\\
\hline
from \textcolor{gray}{Package} import \textcolor{gray}{Methode} & Läd nur einzelne Methoden oder Packete aus einer Bibliothek & from numpy import pi \\
\hline
\end{tabular}
\end{center}

\section*{Arithmetik}

\begin{center}
\begin{tabular}{|p{3cm}|p{8cm}|p{3cm}|}
\hline
a + b & Addition & a + 3 \newline{} \texttt{=>} 5 \\
\hline
a - b & Subtraktion & a - 3.5 \newline{} \texttt{=>} -1.5 \\
\hline
a * b & Multiplikation & a * 5 \newline{} \texttt{=>} 10 \\
\hline
a / b & Division & a / 3 \newline{} \texttt{=>} 0.666666666666666 \\ %Echte Anzahl lassen oder Kürzen?
\hline
a // b & Division ohne Rest & a // 3 \newline{} \texttt{=>} 0 \\
\hline
a \% b & Rest der Division & a \% 3 \newline{} \texttt{=>} 2 \\
\hline
a ** b & Potenzieren ($a^b$) & a ** 3 \newline{} \texttt{=>} 8 \\
\hline
round(x) & Rundet nach Standard Rundungsregeln & round(2.5) \newline{} \texttt{=>} 3 \\
\hline
\end{tabular}
\end{center}

\section*{Vergleichen}

\begin{center}
\begin{tabular}{|p{3cm}|p{8cm}|p{3cm}|}
\hline
a == b & a gleich b & a == 3 \newline{} \texttt{=>} False \\
\hline
a != b & a ungleich b & a != 3 \newline{} \texttt{=>} True \\
\hline
a \texttt{<} b & a kleiner b & a \texttt{<} 3 \newline{} \texttt{=>} True \\
\hline
a \texttt{<=} b & a kleiner oder gleich b & a \texttt{<=} 3 \newline{} \texttt{=>} True \\
\hline
a \texttt{>} b & a größer b & a \texttt{>} 3 \newline{} \texttt{=>} False \\
\hline
a \texttt{>=} b & a größer oder gleich b & a \texttt{>=} 3 \newline{} \texttt{=>} False \\
\hline
\end{tabular}
\end{center}

\section*{Listen}

\begin{center}
\begin{tabular}{|p{3cm}|p{8cm}|p{3cm}|}
\hline
Liste = [x, y, z] & Erzeugt eine Liste mit beliebigen Einträgen & L = [a, 2 * a, 2 * a -1, 1] \newline{} print(L) \newline{} \texttt{=>} [2, 4, 3, 1] \\
\hline
Liste[x] & Gibt den Eintrag der Liste aus & L[a] \newline{} \texttt{=>} 3 \\
\hline
Liste[x:y:z] & Gibt von Eintrag x bis Eintrag y jeden z-ten Eintrag aus. Hierbei ist Liste[x:y] := Liste[x:y:1] & L[:3:2] \newline{} \texttt{=>} [2, 3] \\
\hline
len(Liste) & Anzahl der Listeneinträge & len(L) \newline{} \texttt{=>} 4 \\
\hline
sum(Liste) & Summiert alle Listeneinträge auf & sum(L) \newline{} \texttt{=>} 10 \\
\hline
%Anhängen und entfernen von Elementen hinzufügen? In Datenauswetung wird auf eigene Recherche gesetzt
\end{tabular}
\end{center}

%Methoden werden erstmal ausgelassen, da zu kurz für eigenes Kapitel. Eventuell in Tricks
\newpage
\section*{NumPy}

Für diese Befehle muss zuerst NumPy importiert werden. Hier wird im vollgendem immer np als Abkürzung für NumPy verwendet.

\begin{center}
\begin{tabular}{|p{3cm}|p{8cm}|p{3cm}|}
\hline
np.Methode & Lässt eine Methode aus NumPy ausführen & np.pi \newline{} \texttt{=>} 3.141592653589793 \\
\hline
np.sqrt(x) & Wurzel & sqrt(a) \newline{} \texttt{=>} 1.4142135623730951 \\
\hline
np.array(Liste) & Erstellt aus einer Liste einen NumPy Array & nL = np.Array(L) \\
\hline
array.min() & Gibt das Minimum von einem Array aus & nL.min() \newline{} \texttt{=>} 1 \\
\hline
array.max() & Gibt das Maximum von einem Array aus & nL.max() \newline{} \texttt{=>} 4 \\
\hline
array.sum() & Gibt die Summe von einem Array aus & nL.sum() \newline{} \texttt{=>} 10 \\
\hline
array.mean() & Gibt den Mittelwert von einem Array aus & nL.mean() \newline{} \texttt{=>} 2.5 \\
\hline
array.std() & Gibt die Standartabweichung von einem Array aus & nL.std() \newline{} \texttt{=>} 1.118033988749895 \\
\hline
x, y, z = np.loadtxt( \textcolor{gray}{Dateiname}, \textcolor{gray}{delimiter}, \textcolor{gray}{skiprows}, \textcolor{gray}{usecols}, \textcolor{gray}{unpack}) & Erstellt Arrays x, y und z au den Spalten von der Datei & t, a\_x = np.loadtxt( "Messung1.csv", delimiter=",", skiprows=1, usecols=(0,1), unpack=True) \\ 
&&\\
& Dateiname ist der Name der Datei \newline{} \newline{}
delimeter gibt an, welches Zeichen die Daten trennen \newline{} \newline{}
skiprows gibt an wie viele Spalten bei einem delimiter übersprungen werden müssen \newline{} \newline{}
Wenn nur bestimmte Spalten benötigt werden können diese mit usecols angegeben werden \newline{} \newline{}
unpack sorgt dafür, dass in mehrere array aufgeteilt werden kann
& \\
\hline
\end{tabular}
\end{center}

\section*{MatPlotLib}

Für diese Befehle muss zuerst MatPlotLib importiert werden. Hier wird im vollgendem immer plt als Abkürzung für MatPlotLib verwendet. Eine Einführung gibt es im \href{https://matplotlib.org/stable/users/explain/quick_start.html}{MatPlotLib Quick start Guide}. Hier sollen nur noch einmal eine Auswahl wichtiger Befehle zur schnellen recherche angegeben werden.
Die Ausgaben werden im Folgendem nicht mehr mit angegeben, da diese meist Plots sind. Am besten einfach selber mal eingeben.

\begin{center}
\begin{tabular}{|p{3cm}|p{8cm}|p{3cm}|}
\hline
plt.plot(\textcolor{gray}{x-Achse}, \textcolor{gray}{y-Achse}) & Erstellt einen Plot der Daten, solange die Längen der Listen gleich sind. Die meisten der folgenden Befehle können auch direkt in diesem mit Eingebunden werden. & plt.plot([1, 2, 3], [3, 2, 3]) \\
\hline
plt.xlim(a, b) & Legt den geplotteten x-Bereich fest & plt.xlim(0, 1) \\
\hline
plt.ylim(a, b) & Legt den geplotteten y-Bereich fest & plt.ylim(0, 1) \\
\hline
plt.xlabel() & Fügt eine x-Achsenbeschriftung hinzu & plt.xlabel('Zeit in s') \\
\hline
plt.ylabel() & Fügt eine y-Achsenbeschriftung hinzu & plt.xlabel('Strecke in m') \\
\hline
\end{tabular}
\end{center}

\section*{Tricks} %Soll am beim fertigstellung des Glossars ans Ende kommen

\begin{center}
\begin{tabular}{|p{3cm}|p{8cm}|p{3cm}|}
\hline
str(x) & Erstellt einen String aus x sofern möglich & string\_a = str(a) \\
\hline
int(x) & Erstellt einen Integer aus x sofern möglich & int(string\_a) \\
\hline
x. & Erstellt statt einem Integer einen Float & 
\begin{tabular}{|p{1.2cm}|p{1.1cm}|}
2 & 2. \\
\texttt{=>} 2 & \texttt{=>} 2.0 \\
\end{tabular}
\\
\hline
\end{tabular}
\end{center}

%Standart Tabelle
\begin{comment}
\begin{center}
\begin{tabular}{|p{3cm}|p{8cm}|p{3cm}|}
\hline
 & & \\
\hline
\end{tabular}
\end{center}
\end{comment}


\end{document}